\chapter{Background}\label{ch:background}

This chapter collects the mathematical facts used later. Notation follows
Chapter~\ref{ch:intro}, and the definitions introduced here are referred to
from Chapter~\ref{ch:model}.

\section{Probability preliminaries}\label{sec:prob}

Let $(\Omega,\mathcal{F},\mathbb{P})$ be a probability space. A random vector
$X$ taking values in $\mathbb{R}^d$ has mean $\mu = \mathbb{E}[X]$ and
covariance matrix $\Sigma = \mathbb{E}[(X-\mu)(X-\mu)^\top]$, which is
symmetric and positive semidefinite.

\lipsum[7-8]

\section{Linear algebra}\label{sec:linalg}

For a symmetric matrix $A$ with eigendecomposition $A = Q\Lambda Q^\top$, the
matrix function $f(A)$ is defined by
\begin{equation}\label{eq:matfun}
  f(A) = Q\,\operatorname{diag}\bigl(f(\lambda_1),\dots,f(\lambda_d)\bigr)\,Q^\top .
\end{equation}
Equation~\eqref{eq:matfun} is used in Section~\ref{sec:estimators}.

\begin{table}[ht]
  \centering
  \begin{tabular}{lll}
    \toprule
    Symbol & Meaning & Typical size \\
    \midrule
    $n$ & number of observations & $10^2$--$10^4$ \\
    $d$ & dimension of each observation & $2$--$50$ \\
    $k$ & number of mixture components & $1$--$8$ \\
    \bottomrule
  \end{tabular}
  \caption{Symbols used throughout the report.}\label{tab:symbols}
\end{table}

Table~\ref{tab:symbols} lists the recurring symbols and their typical sizes.

\lipsum[9-10]

\section{Further remarks}\label{sec:further-background}

The remarks in this section restate the main point of the chapter in more
general terms, so that the argument can be read without the surrounding
derivations.

\lipsum[63-74]
